\documentclass[11pt]{article}
\usepackage[a4paper,margin=2.05cm]{geometry}
\usepackage[T1]{fontenc}
\usepackage{lmodern}
\usepackage{amsmath,amssymb}
\usepackage{xcolor}
\usepackage{microtype}
\usepackage[hidelinks]{hyperref}
\setlength{\emergencystretch}{3em}

\definecolor{ink}{HTML}{172033}
\definecolor{accent}{HTML}{385B7A}
\definecolor{soft}{HTML}{EEF5FA}
\color{ink}

\title{The two-matrix formula for a periodic tridiagonal numerical range}
\author{Literature answer to Conjecture 3.7 in arXiv:1910.00720}
\date{17 August 2026}

\begin{document}
\maketitle

\begin{center}
\fcolorbox{accent}{soft}{\parbox{0.91\textwidth}{
\textbf{Status: literature-already-answered, full conjecture.}
For every $n$, the closure of the numerical range of the periodic
tridiagonal operator in Conjecture~3.7 is the convex hull of the numerical
ranges of the two stated $(n+1)\times(n+1)$ matrices.  The later paper
arXiv:2103.01866 explicitly proves this conjecture.}}
\end{center}

\section*{Source conjecture}
Let $b$ be the $(n+1)$-periodic sequence with period word $0^n1$, and let
$T=T(b,0,1)$ be the associated tridiagonal operator.  Let $B_n$ be the
$(n+1)\times(n+1)$ matrix with ones on the first superdiagonal and zeros
elsewhere, and let
\[
 J_n=\operatorname{diag}(1,0,\ldots,0,1).
\]
Conjecture~3.7 of Itzá-Ortiz--Martínez-Avendaño, on source PDF page~19,
asserts that
\begin{equation}\label{eq:source}
 \overline{W(T)}=\operatorname{co}\!\left(
 W(B_n+J_n)\cup W(B_n-J_n)\right).
\end{equation}
The source proves the case $n=1$ and reports numerical verification for
$n=2,3$, but says that it cannot yet provide a proof \cite{IM}.  Its
introduction also records that the referee had proposed a proof expected to
appear later.

\section*{Explicit later solution}
The abstract of Itzá-Ortiz--Martínez-Avendaño--Nakazato states that the paper
proves a conjecture of the first two authors, and its introduction identifies
that conjecture as Conjecture~3.7 of the source \cite{IMN}.

Their Theorem~2.3, beginning on supporting PDF page~7, proves a more general
statement.  For suitable real $(n+1)$-periodic lower- and upper-diagonal words
$a,c$ satisfying a palindromic system of modulus equalities, it constructs two
explicit tridiagonal matrices $A^+$ and $A^-$ and proves
\[
 \overline{W(T(a,0,c))}
 =\operatorname{co}\!\left(W(A^+)\cup W(A^-)\right).
\]

Example~2.4, on supporting PDF pages~11--12, takes $a$ with period word
$0100\cdots0$ and $c$ with period word $11\cdots1$.  The hypotheses of
Theorem~2.3 hold, and its two matrices are exactly $B_n+J_n$ and $B_n-J_n$.
After passing by unitary equivalence to the period word $0^n1$ used in the
source, the example concludes explicitly that it proves Conjecture~3.7.
Consequently, \eqref{eq:source} holds for every $n$.

\section*{Mechanism and scope}
The proof encodes every support function of the infinite operator's numerical
range in a homogeneous Kippenhahn polynomial.  Under the palindromic
hypotheses, this polynomial factors into the Kippenhahn polynomials of
$A^+$ and $A^-$.  The block-diagonal matrix $A^+\oplus A^-$ therefore has the
same closed numerical range, and its numerical range is the convex hull of the
two block numerical ranges.

The later result is stronger than the source conjecture because it handles a
larger palindromic family.  The exact specialization leaves no part of
Conjecture~3.7 open.  This packet records an explicit literature answer, not an
agent-generated proof.

\section*{Evidence checked}
The source and supporting arXiv sources were rebuilt locally.  The decisive
locations are source PDF page~19, supporting Theorem~2.3 on page~7, and
supporting Example~2.4 on pages~11--12.  The arXiv records and journal DOI
metadata were also checked.

{\raggedright
\begin{thebibliography}{9}
\bibitem{IM}
B. A. Itzá-Ortiz and R. A. Martínez-Avendaño,
\emph{The numerical range of a class of periodic tridiagonal operators},
arXiv:1910.00720; Linear Multilinear Algebra \textbf{69} (2021), 786--806,
\href{https://doi.org/10.1080/03081087.2019.1706438}{doi:10.1080/03081087.2019.1706438}.

\bibitem{IMN}
B. A. Itzá-Ortiz, R. A. Martínez-Avendaño, and H. Nakazato,
\emph{The numerical range of some periodic tridiagonal operators is the
convex hull of the numerical ranges of two finite matrices},
arXiv:2103.01866; Linear Multilinear Algebra \textbf{69} (2021), 2830--2849,
\href{https://doi.org/10.1080/03081087.2021.1957760}{doi:10.1080/03081087.2021.1957760}.
\end{thebibliography}
\par}

\end{document}
